\par\medskip\noindent\textbf{The exact two-boundary core.}
On a Hopf core the two distinct marked curves have the actual
reduced anticanonical divisor
$-K_H=J_-+J_+$. Choose a finite Galois primary cover:
a normal infinite cyclic subgroup is obtained by intersecting
the finitely many conjugates of a finite-index primary subgroup.
The inverse images of the two marked curves are disjoint reduced
divisors. The primary surface cannot be nonlinear, whose only
compact curve is its unique axis. In diagonal coordinates its
contraction is $F(z,w)=(\alpha z,\beta w)$.

Pull back the anticanonical section. In the frame
$\partial_z\wedge\partial_w$ its coefficient $g$ is holomorphic
on $\C^2\setminus0$ and extends across zero by Hartogs. The exact
deck identity is
\[
                    g(\alpha z,\beta w)=\alpha\beta\,g(z,w).
\]
For a monomial $z^a w^b$ with $a,b\geq1$, comparison of
absolute values forces $a=b=1$. A pure power $w^{r+1}$ can
occur only when $\alpha=\beta^r$, and a pure power $z^{s+1}$
only when $\beta=\alpha^s$. Both occur only in the scalar
case $\alpha=\beta$, $r=s=1$. Consequently in a nonscalar
resonant case the reduced divisor is
$g=cw(z+aw^r)$ with $c\ne0$, or its coordinate interchange;
in the nonresonant case it is $g=czw$; and in the scalar
case it is a quadratic with two distinct linear factors.
The triangular map $z\mapsto z+aw^r$, or the scalar linear
change of basis, commutes with $F$ and sends this exact
divisor to the two axes.

Each marked curve therefore has just one connected preimage
axis in the primary cover. Every finite deck lift preserves
both axes. It extends with its inverse across the origin,
fixes the origin, and conjugates $F$ to $F^{\pm1}$.
Conjugation to $F^{-1}$ is impossible already on tangent
spaces: the eigenvalue moduli of $DF(0)$ are less than one,
whereas those of its inverse are greater than one.
The lift thus commutes with $F$. Write it as
$(z\,a(z,w),w\,b(z,w))$. Equivariance gives
$a\circ F=a$ and $b\circ F=b$; contraction to the origin
makes both functions constant. Hence the full quotient
group is diagonal and the two marked curves are exactly
its distinct axes. This also treats the initially possible
nonabelian scalar secondary cover: preservation of these
two reduced boundary curves forces its actual lifted group
to be diagonal.

On a minimal ruled core let $p:R\to B$ be its ruling over
the elliptic curve, and retain its actual elliptic
quasi-bundle $q:R\to\PP^1$ and the two marked fibre supports
$J_-,J_+$. A ruling fibre $f$ has $K_R\cdot f=-2$.
The identity $K_R=-J_--J_+$ gives
\[
                 \deg(p|_{J_-})+\deg(p|_{J_+})=2.
\]
Both degrees are positive: an elliptic curve cannot lie
in a rational ruling fibre. Each degree is therefore one,
and the actual maps $J_\pm\to B$ are isomorphisms.
The two disjoint sections give two quotient lines of the
rank-two ruling bundle. Their combined map is an
isomorphism on each fibre, since their two quotient
kernels are distinct. It is consequently a bundle
isomorphism, giving the marked surface
$R=\PP(\OO_B\oplus L)$.

The retained map $q$, not merely the canonical identity,
determines the degree and order of $L$. Write
$q^*(0)=m_-J_-$ and $q^*(\infty)=m_+J_+$.
Each full fibre has square zero, so $J_\pm^2=0$ and
$\deg L=0$. Intersecting their linear equivalence with
$f$ gives $m_-=m_+=m$. Their actual divisor difference
trivializes $L^m$. If $L$ had order $d<m$, its fibre
coordinate would give a meromorphic function $q_0$
with divisor $d(J_--J_+)$, and
$q=cq_0^{m/d}$. This contradicts connectedness of the
regular fibres of $q$. Hence $L$ has exact order $m$.
On its connected finite \'{e}tale trivializing cover
$B'\to B$, the actual model is $B'\times\PP^1$ with
a primitive cyclic translation on $B'$ and the inverse
primitive fibre rotation; its quotient map is $q=z^m$.
This identifies the two marked full-fibre orders and
the regular elliptic Jacobian, as well as the ruled
surface and its normal lines.
